\documentclass[11pt]{article}
\usepackage[margin=1in]{geometry}
\usepackage{booktabs}
\usepackage{longtable}
\usepackage{hyperref}
\usepackage{listings}
\usepackage[T1]{fontenc}

\title{Sparse Blackwell Tensor Core Microbenchmark Plan and Report}
\author{microbench-blackwell}
\date{May 9, 2026}

\begin{document}
\maketitle

\section{Scope}

This report defines benchmark coverage for sparse Blackwell tensor-core MMA
instructions and their dense counterparts. The executable benchmark in this
directory targets the fifth-generation Tensor Core PTX instructions
\texttt{tcgen05.mma}, \texttt{tcgen05.mma.sp}, \texttt{tcgen05.mma.ws}, and
\texttt{tcgen05.mma.ws.sp}.

The sparse comparison is instruction-pair based: \texttt{tcgen05.mma.sp} is
compared with \texttt{tcgen05.mma}, and \texttt{tcgen05.mma.ws.sp} is compared
with \texttt{tcgen05.mma.ws}. The existing \texttt{umma\_throughput/} and
\texttt{umma\_latency/} benchmarks already cover a subset of dense
\texttt{tcgen05.mma} behavior; this directory keeps sparse work isolated while
using the same throughput and latency measurement model.

Implemented artifacts are \texttt{tcgen05\_sparse\_mma.cu}, the
\texttt{benchmark.py} sweep driver, a named-format \texttt{Makefile}, and this
Markdown/TeX/PDF report set.

\section{Instruction Families}

\begin{longtable}{llll}
\toprule
Family & Sparse PTX kind & Dense peer & Data formats covered \\
\midrule
\endhead
Standard MMA & \texttt{tf32} & \texttt{tf32} & TF32 \\
Standard MMA & \texttt{f16} & \texttt{f16} & FP16, BF16 \\
Standard MMA & \texttt{f8f6f4} & \texttt{f8f6f4} & E4M3, E5M2, E2M3, E3M2, E2M1 \\
Standard MMA & \texttt{i8} & \texttt{i8} & S8, U8 \\
Standard MMA, block scaled & \texttt{mxf8f6f4.block\_scale} & same & MX narrow formats \\
Standard MMA, block scaled & \texttt{mxf4.block\_scale} & same & MXFP4 \\
Standard MMA, block scaled & \texttt{mxf4nvf4.block\_scale} & same & MXFP4, NVFP4 \\
Weight-stationary MMA & \texttt{tf32} & \texttt{tf32} & TF32 \\
Weight-stationary MMA & \texttt{f16} & \texttt{f16} & FP16, BF16 \\
Weight-stationary MMA & \texttt{f8f6f4} & \texttt{f8f6f4} & E4M3, E5M2, E2M3, E3M2, E2M1 \\
Weight-stationary MMA & \texttt{i8} & \texttt{i8} & S8, U8 \\
\bottomrule
\end{longtable}

\texttt{tcgen05.mma.ws.sp} is CTA-group 1 only. Standard
\texttt{tcgen05.mma.sp} supports CTA groups 1 and 2. The benchmark driver
skips unsupported combinations rather than failing late.

\section{Sparse Semantics}

Sparse \texttt{tcgen05.mma.sp} and \texttt{tcgen05.mma.ws.sp} use sparse matrix
\texttt{A}. The benchmark sets the PTX instruction descriptor sparsity bit and
stores sparsity metadata in Tensor Memory.

\begin{tabular}{lll}
\toprule
Kind & Sparse granularity & Metadata pattern \\
\midrule
\texttt{tf32} & 1:2 & repeated \texttt{0b0100} nibbles \\
\texttt{f16} & 2:4 & repeated \texttt{0b0100} nibbles \\
\texttt{f8f6f4} & 2:4 & repeated \texttt{0b0100}, selector 0 \\
\texttt{mxf8f6f4} & 2:4 & repeated \texttt{0b0100}, selector 0 \\
\texttt{i8} & 2:4 & repeated \texttt{0b0100}, selector 0 \\
\texttt{mxf4} & 4:8 in pairs & repeated \texttt{0b0100}, selector 0 \\
\texttt{mxf4nvf4} & 4:8 in pairs & repeated \texttt{0b0100}, selector 0 \\
\bottomrule
\end{tabular}

\section{Measurement Method}

Throughput mode issues a burst of \texttt{MMA\_DEPTH} operations in each
iteration and reports cycles per MMA:

\begin{lstlisting}
cycles_per_mma = total_clock_cycles / (MMA_DEPTH * NUM_ITERS)
\end{lstlisting}

Latency mode issues one MMA per sample, commits it to an mbarrier, waits for
completion, sorts the samples, and reports the median cycle count.

The benchmark reports both raw instruction throughput and dense-equivalent
operations. Sparse instructions store half of matrix \texttt{A}, but their
dense-equivalent GEMM shape remains \texttt{M x N x K}. CSV output therefore
keeps both \texttt{MathOpsPerMMA = 2*M*N*K} and \texttt{StoredAElements}.

\section{Usage}

\begin{lstlisting}[language=bash]
cd tcgen05_sparse_mma
make clean
make MODE=throughput INSTRUCTION=mma KIND=f16 A_TYPE=f16 B_TYPE=f16 \
     SPARSITY=1 MMA_M=128 MMA_N=64 MMA_K=32 MMA_DEPTH=256 \
     CTA_GROUP=1 AB_LAYOUT=ss
./tcgen05_sparse_mma.out
\end{lstlisting}

\begin{lstlisting}[language=bash]
python3 benchmark.py --mode throughput --preset smoke --sparsity both \
    --output sparse_tput.csv --overwrite
python3 benchmark.py --mode latency --preset smoke --sparsity both \
    --output sparse_lat.csv --overwrite
\end{lstlisting}

Build without running, useful on machines that have a CUDA toolchain but not an
SM100/SM110 GPU:

\begin{lstlisting}[language=bash]
python3 benchmark.py --mode throughput --preset smoke --sparsity both \
    --instruction all --ab-layout ss --cta-group 1 --build-only \
    --output sparse_build_matrix.csv --overwrite
\end{lstlisting}

The default \texttt{CUDA\_GENCODE} is:

\begin{lstlisting}
-gencode arch=compute_100a,code=sm_100a
\end{lstlisting}

This explicit \texttt{a} target is required. On the local CUDA 13.1 install,
\texttt{-arch=sm\_100a} lowered to plain \texttt{sm\_100} during ptxas, which
rejects \texttt{tcgen05.*}.

\section{Platform Notes}

The code targets Blackwell server-family tensor-core instructions. The local
RTX 5090 environment documented elsewhere in this repository reports
\texttt{sm\_120}, and that target rejects \texttt{tcgen05.*} instructions. Full
runtime validation requires an SM100/SM110-family target such as
\texttt{sm\_100a} or \texttt{sm\_110a}.

Local validation completed: Python syntax checking for \texttt{benchmark.py},
TeX build to \texttt{REPORT.pdf}, and representative \texttt{sm\_100a}
executable builds for dense \texttt{tcgen05.mma.kind::f16}, sparse
\texttt{tcgen05.mma.sp.kind::f16}, the sparse TS operand form, sparse
\texttt{tcgen05.mma.ws.sp.kind::f16}, sparse block-scaled \texttt{mxf8f6f4},
\texttt{mxf4}, and \texttt{mxf4nvf4}, and sparse CTA-group 2 \texttt{i8}. Local
runtime produces the expected launch error on this host:
\texttt{no kernel image is available for execution on the device}.

\section{References}

\begin{itemize}
\item NVIDIA PTX ISA 9.1, sparse \texttt{tcgen05.mma.sp} syntax and semantics:
\url{https://docs.nvidia.com/cuda/archive/13.1.1/parallel-thread-execution/index.html#tensorcore-5th-generation-instructions-tcgen05-mma-sp}
\item NVIDIA PTX ISA 9.1, sparse matrix formats:
\url{https://docs.nvidia.com/cuda/archive/13.1.1/parallel-thread-execution/index.html#sparse-matrices}
\item NVIDIA PTX ISA 9.1, instruction descriptor:
\url{https://docs.nvidia.com/cuda/archive/13.1.1/parallel-thread-execution/index.html#instruction-descriptor}
\item NVIDIA CUTLASS Blackwell functionality overview:
\url{https://docs.nvidia.com/cutlass/latest/media/docs/cpp/blackwell_functionality.html}
\end{itemize}

\end{document}
